\section*{Capítulo \ref{ch:b3:supramolecular} --- Química supramolecular}

\begin{solution}{exo:b3:supramolecular:1}
Ion--dipolo; apilamiento $\pi$; catión--$\pi$; \omterm{def:b3:supramolecular:interactions}{enlace de halógeno} (de I a N); tres enlaces de hidrógeno.
\end{solution}

\begin{solution}{exo:b3:supramolecular:2}
$\Delta G^\circ = -RT\ln K = -8.314 \times 298 \times \ln(\num{1.0e4}) = \qty{-22.8}{kJ/mol}$.
\end{solution}

\begin{solution}{exo:b3:supramolecular:3}
$K_{\ce{K+}}/K_{\ce{Na+}} = 10^{1.7} = 50$; la diferencia es de $RT\ln 50 = \qty{9.7}{kJ/mol}$ a favor del potasio.
\end{solution}

\begin{solution}{exo:b3:supramolecular:4}
$n/(n + m) = 1/3$: un \omterm{def:b3:supramolecular:supramolecular}{anfitrión} por cada dos \omterm{def:b3:supramolecular:supramolecular}{huéspedes}, $\mathrm{HG_2}$.
\end{solution}

\begin{solution}{exo:b3:supramolecular:5}
$H_0 + G_0 + 1/K = \qty{4.667e-3}{mol/L}$; $[\mathrm{HG}] = (4.667 - \sqrt{4.667^2 - 4 \times 2.0 \times 2.0})/2 \times 10^{-3} = \qty{1.13e-3}{mol/L}$: fracción
unida, 0.566; $\delta = 3.560 + 0.160 \times 0.566 = \qty{3.651}{ppm}$.
\end{solution}

\begin{solution}{exo:b3:supramolecular:6}
$\log K = 18.3 - 8.6 = 9.7$, $K = \num{5e9}$; $\Delta G^\circ = -RT\ln K = \qty{-55}{kJ/mol}$. Los seis enlaces Ni--N son del
mismo tipo a ambos lados; el número de moléculas libres pasa de cuatro a
siete: la ganancia de entropía de traslación impulsa la reacción.
\end{solution}

\begin{solution}{exo:b3:supramolecular:7}
$\Delta G^\circ = -RT\ln(\num{2.0e5}) = \qty{-30.2}{kJ/mol}$; $T\Delta S^\circ = \Delta H^\circ - \Delta G^\circ = -40 + 30.2 = \qty{-9.8}{kJ/mol}$;
$\Delta S^\circ = \qty{-33}{J.K^{-1}.mol^{-1}}$: una unión impulsada por la entalpía, pagada en parte con entropía.
\end{solution}

\begin{solution}{exo:b3:supramolecular:8}
El cobre(I) une dos unidades bidentadas de fenantrolina y, al ser tetraédrico, las mantiene
en ángulo recto, una enhebrada a través del espacio en el que se cerrará el anillo de la
otra. Cerrar después ambas cadenas en anillos las entrelaza. Un gran exceso de
cianuro, que se une al cobre(I) con mucha más fuerza, elimina el metal y deja
el \omterm{def:b3:supramolecular:interlocked}{catenano} libre.
\end{solution}

\begin{solution}{exo:b3:supramolecular:9}
$k_c = \pi\Delta\nu/\sqrt2 = \pi \times 120/1.414 = \qty{267}{s^{-1}}$. Eyring:
$\Delta G^\ddagger = RT\ln(k_BT/hk_c) = 8.314 \times 300 \times \ln(\num{6.25e12}/267) = \qty{59.6}{kJ/mol}$.
\end{solution}

\begin{solution}{exo:b3:supramolecular:10}
Con un exceso de fármaco sólido, el fármaco libre queda fijado en $S_0$. Entonces $[\mathrm{DC}] = KS_0[\mathrm C]$ y
$C_t = [\mathrm C] + [\mathrm{DC}] = [\mathrm C](1 + KS_0)$, luego $[\mathrm{DC}] = KS_0C_t/(1 + KS_0)$, y la solubilidad total es $S_0 + [\mathrm{DC}]$. Aquí,
$KS_0 = 0.050$: $S = 0.10 + 0.050 \times 10/1.050 = \qty{0.58}{mmol/L}$, casi seis veces el valor intrínseco.
\end{solution}

\begin{solution}{exo:b3:supramolecular:11}
Cuando $1/K \ll H_0, G_0$, $[\mathrm{HG}] \approx \bigl(H_0 + G_0 - |H_0 - G_0|\bigr)/2 = \min(H_0, G_0)$: la fracción unida es $G_0/H_0$ hasta un
equivalente, y después 1. La curva ya no depende de $K$: cualquier valor lo bastante grande
da el mismo quiebro neto dentro del ruido, de modo que los datos solo muestran que $K$ supera
cierto mínimo.
\end{solution}

\begin{solution}{exo:b3:supramolecular:12}
$k_{\mathrm{off}} = k_{\mathrm{on}}/K = \qty{1e9}{L.mol^{-1}.s^{-1}}/\qty{1500}{L.mol^{-1}} = \qty{6.7e5}{s^{-1}}$. La diferencia de desplazamiento,
$\qty{0.160}{ppm}$, equivale a \qty{64}{Hz} a \qty{400}{MHz}; la coalescencia exigiría $k \approx \pi \times 64/\sqrt2 = \qty{140}{s^{-1}}$. El intercambio es
miles de veces más rápido: una sola señal promediada.
\end{solution}

\begin{solution}{pb:b3:supramolecular:1}
\textbf{1.} Un anillo de seis unidades $\ce{-CH2CH2O-}$; en el complejo, el $\ce{K+}$ se sitúa en el centro, unido a los
seis oxígenos, que apuntan hacia dentro, con los grupos $\ce{CH2}$ en el exterior.

\textbf{2.} $\ce{K+}$: $-RT\ln 10 \times 6.1 = \qty{-34.8}{kJ/mol}$; $\ce{Na+}$: \qty{-25.1}{kJ/mol}.

\textbf{3.} $10^{6.1 - 4.4} = 50$.

\textbf{4.} El $\ce{K+}$, más grande (\qty{138}{pm}), alcanza a la vez los seis oxígenos del anillo; el
$\ce{Na+}$, más pequeño (\qty{102}{pm}), no puede, salvo que el anillo se pliegue, lo que cuesta tensión.

\textbf{5.} Su exterior es una superficie hidrocarbonada y su carga está enterrada; el $\ce{KCl}$
necesitaría que sus iones estuvieran solvatados, lo que el tolueno no puede hacer.

\textbf{6.} El agua solvata el $\ce{K+}$ mucho mejor que el metanol, y el ion debe perder esa
agua para entrar en el anillo.

\textbf{7.} $\ce{K+}(\text{aq}) + \ce{MnO4-}(\text{aq}) + \mathrm L(\text{org}) \rightleftharpoons \ce{[KL]+MnO4-}(\text{org})$.

\textbf{8.} $[\ce{K+}]_{\mathrm{aq}}$ se mantiene en \qty{0.10}{mol/L}: $y = \num{1.0e5} \times 0.10 \times (\num{1.0e-3} - y)^2 = \num{1.0e4}(\num{1.0e-3} - y)^2$.

\textbf{9.} Con $u = \num{1.0e-3} - y$: $\num{1.0e4}u^2 + u - \num{1.0e-3} = 0$, $u = \qty{2.70e-4}{mol/L}$, $y = \qty{7.30e-4}{mol/L}$: se extrae
el 73\,\%.

\textbf{10.} $y = \num{1.0e4}(\num{1.0e-3} - y)(\num{1.0e-2} - y)$ da $y = \qty{9.89e-4}{mol/L}$: el 99\,\%.

\textbf{11.} $y = \num{1.0e4}(\num{1.0e-3} - y)(\num{1.0e-4} - y)$ da $y = \qty{9.0e-5}{mol/L}$: el 9\,\%, limitado por la corona,
que está cargada al 90\,\%.

\textbf{12.} El ion permanganato conserva su color púrpura en el tolueno. Allí
no está solvatado, es un anión desnudo y reacciona mucho más deprisa que en agua, donde una
capa de hidratación lo protege.

\textbf{13.} El ion sodio entra y sale mucho más deprisa que la diferencia de las
dos frecuencias de resonancia: intercambio rápido.

\textbf{14.} $\delta = \delta_{\mathrm{libre}} + (\delta_{\mathrm{unido}} - \delta_{\mathrm{libre}})[\mathrm{HG}]/H_0$, con $[\mathrm{HG}]$ dada por la isoterma exacta, $H_0 = \qty{2.0e-3}{mol/L}$,
$G_0 = (\text{equiv.})\,H_0$, $\delta_{\mathrm{libre}} = \qty{3.560}{ppm}$.

\textbf{15.} $K = (1.55 \pm 0.08) \times 10^3\ \unit{L.mol^{-1}}$, $\delta_{\mathrm{unido}} = \qty{3.719 \pm 0.001}{ppm}$.

\textbf{16.} El residuo cuadrático medio es de \qty{0.0014}{ppm}, del orden de la precisión
de lectura, sin tendencia: el modelo 1:1 se ajusta.

\textbf{17.} Fracción unida $f$ en $G_0 = fH_0 + f/(K(1 - f))$: 0.28 equivalentes para $f = 0.2$ y 2.1 para $f = 0.8$.

\textbf{18.} $KH_0 = 10^{6.1} \times \num{2.0e-3} \approx 2500$: la curva presentaría un quiebro neto en un equivalente
y solo daría una cota inferior; hace falta un experimento de competencia o una calorimetría a
menor concentración.

\textbf{19.} En el tolueno, en una disolución homogénea de alqueno y permanganato.

\textbf{20.} El manganeso reducido abandona la fase orgánica (como dióxido de manganeso),
y la corona, liberada en la interfase, capta otro $\ce{K+}$ y otro $\ce{MnO4-}$. Cada molécula de
corona transporta permanganato muchas veces: basta una cantidad catalítica.

\textbf{21.} $K_{\mathrm{ex}}$ multiplica $[\ce{K+}]$: una concentración alta de potasio favorece la extracción.

\textbf{22.} Las sales de amonio cuaternario, como el cloruro de tetrabutilamonio o el de
benciltrietilamonio, cuyo catión es lipófilo por sí mismo.

\textbf{23.} El benceno provoca leucemia; el tolueno hace el mismo trabajo con un riesgo
mucho menor.

\textbf{24.} Con \qty{1.0}{mmol/L} de corona, el 73\,\% del permanganato se extrae al
tolueno.
\end{solution}
